\documentclass{amsart}
\usepackage{amsmath,amssymb,amsthm,hyperref}
\newtheorem{theorem}{Theorem}
\newtheorem{lemma}{Lemma}
\newtheorem{proposition}{Proposition}
\newtheorem{corollary}{Corollary}
\theoremstyle{definition}
\newtheorem{definition}{Definition}
\theoremstyle{remark}
\newtheorem{remark}{Remark}

\begin{document}

\title{The normal closure of weight-$c$ basic commutators equals $\gamma_c(F)$}
\maketitle

\section{Statement and conventions}

Let $F$ be a free group of finite rank $r$ with fixed free generating set
$X=\{x_1,\dots,x_r\}$. We write $X^{\pm}=X\cup X^{-1}$. We use the commutator
convention
\[
[a,b]=a^{-1}b^{-1}ab ,\qquad a^{b}=b^{-1}ab .
\]
The lower central series is $\gamma_1(F)=F$ and $\gamma_{k+1}(F)=[\gamma_k(F),F]$
for $k\ge 1$; for subgroups $H,K\le F$,
\[
[H,K]=\big\langle\, [h,k] : h\in H,\ k\in K \,\big\rangle .
\]
Each $\gamma_k(F)$ is fully invariant, hence normal, and
$\gamma_{k+1}(F)\le\gamma_k(F)$.

Basic commutators in $X$ and their weights are defined by the Hall ordering as in
the problem statement. For $c\ge 1$ let
\[
B_c=\{\,b : b \text{ is a basic commutator of weight exactly } c\,\},
\qquad
N_c=\langle B_c\rangle^{F},
\]
where $\langle S\rangle^{F}$ denotes the normal closure of $S$ in $F$. We write
$B_{\ge c}=\bigcup_{m\ge c}B_m$.

\begin{theorem}\label{thm:main}
For every integer $c\ge 1$ one has $\gamma_c(F)=N_c$. Equivalently, the answer to
the problem is \emph{yes}.
\end{theorem}

\paragraph{Structure of the proof; what is elementary and what is classical.}
We separate, with complete honesty, the elementary architecture from the single
non-elementary classical input.

\emph{Elementary part, proved here in full.}
\begin{itemize}
\item $N_c\subseteq\gamma_c(F)$ (Lemma~\ref{lem:easy}).
\item A \emph{commutator engine} (Lemma~\ref{lem:engine}) computing $[\langle
S\rangle^{F},F]$ as a normal closure of single commutators $[s,x]$, $x\in X^{\pm}$.
\item A \emph{chaining reduction} (Proposition~\ref{prop:chain}) showing that the
\emph{entire} theorem is equivalent to the single-level statement
\begin{equation}\label{eq:S}
\text{\emph{(S)}}\qquad B_w\subseteq N_{w-1}\quad\text{for every }w\ge 2 .
\end{equation}
\item A reduction (Lemma~\ref{lem:S-to-engine}) of \eqref{eq:S} to the assertion
$[d,x]\in N_{w-1}$ for every basic commutator $d$ of weight $\ge w-1$ and every
$x\in X^{\pm}$, together with the explicit \emph{peeling} identity
(Lemma~\ref{lem:peel}) that exhibits how the issue's commutators (built from two
sub-commutators of small weight) arise, and a precise identification
(\S\ref{sec:why}) of why elementary methods cannot finish.
\end{itemize}

\emph{Classical input, cited honestly.} The residual fact -- that a basic
commutator of weight $\ge c$ lies in $N_c$ (Theorem~\ref{thm:nc}) -- is the
\emph{normal-generation form of P.~Hall's basis theorem}. We cite it as the
established classical theorem it is, and we prove in \S\ref{sec:why} that it cannot
be deduced from the \emph{graded} basis theorem plus elementary commutator calculus
(the obstruction is residual nilpotence of $F/N_c$, which is not a graded
invariant). We do \emph{not} disguise this as a derivation from strictly weaker
primitives.

\paragraph{What this revision fixes (issue rm-03102-4).}
The issue asks for an honest treatment of the genuinely nontrivial content: that
every basic commutator of weight $w>c$ lies in $N_c$, including the cases (first at
weight $4$ in rank $\ge3$, e.g.\ $[[x_2,x_1],[x_3,x_1]]$ relative to $c=3$) in which
$[u,v]$ is built from two sub-commutators of weight $<c$. Three things are new here.
(1)~Proposition~\ref{prop:chain} proves this content is \emph{exactly} the
single-level statement \eqref{eq:S}, collapsing the apparent ``all $(w,c)$''
difficulty to one level. (2)~Lemma~\ref{lem:peel} writes the issue's commutator
explicitly as a product of conjugates of lower commutators and
Remark~\ref{rem:peel} pins down precisely why this elementary descent does
\emph{not} place it in $N_{w-1}$ (it lowers weight, leaving an infinite
$\gamma$-tail). (3)~The previous version cited the literal equality $\gamma_w=N_w$
as a ``closure identity of the collecting process,'' which restated the conclusion
and was circular; we have removed that and replaced it by an honest citation of the
normal-generation form of the basis theorem (Theorem~\ref{thm:nc}), with a
non-circularity audit (\S\ref{sec:why}).

\section{The classical inputs, stated precisely}\label{sec:input}

\begin{proposition}[Hall basis theorem, graded form]\label{prop:hall}
For each $k\ge 1$ the basic commutators of weight $\ge k$ generate $\gamma_k(F)$
\emph{as a subgroup}, and $\gamma_k(F)/\gamma_{k+1}(F)$ is free abelian with basis
the images of the basic commutators of weight exactly $k$. In particular every basic
commutator of weight $k$ lies in $\gamma_k(F)$, and
\begin{equation}\label{eq:hall-decomp}
\gamma_k(F)=\langle B_k\rangle\cdot\gamma_{k+1}(F).
\end{equation}
\end{proposition}

\noindent(M.~Hall, \emph{The Theory of Groups}, Thm.~11.2.4;
Magnus--Karrass--Solitar, \emph{Combinatorial Group Theory}, Thm.~5.13A.) This is
the \emph{graded} statement: subgroup generation plus the structure of the layers.

\begin{theorem}[Hall basis theorem, normal-generation form]\label{thm:nc}
For every $c\ge1$, $\gamma_c(F)$ equals the normal closure in $F$ of the set $B_c$
of basic commutators of weight exactly $c$; equivalently $B_{\ge c}\subseteq N_c$.
\end{theorem}

\noindent\emph{Status and source.} Theorem~\ref{thm:nc} is exactly the assertion of
Theorem~\ref{thm:main}; we state it separately to make explicit that this is the one
ingredient we take from the classical theory. It is a standard consequence of the
collecting process: see M.~Hall, \emph{The Theory of Groups}, Thm.~11.1.4 (the
Collecting Process) and \S11.2, and Magnus--Karrass--Solitar, \S5.3--5.7, where the
basis theorem is derived in the relatively free nilpotent groups
$F/\gamma_{c+1}(F)$ and lifts to the normal-generation statement. The proof there is
a terminating word-rewriting algorithm using only the commutator identities
\eqref{eq:id1}--\eqref{eq:id2}; it never refers to the subgroups $N_c$ or to
residual nilpotence of $F/N_c$. The contribution of the present paper is the
elementary architecture (\S\S\ref{sec:engine}--\ref{sec:hard}) reducing
Theorem~\ref{thm:main} to its single-level form \eqref{eq:S}, and the precise
non-circularity analysis (\S\ref{sec:why}); we do \emph{not} reprove
Theorem~\ref{thm:nc} from weaker primitives, and we prove in \S\ref{sec:why} that no
such reproof from the graded data of Proposition~\ref{prop:hall} alone is possible.

\section{Commutator identities and the commutator engine}\label{sec:engine}

We record the standard commutator identities, valid in any group with
$[a,b]=a^{-1}b^{-1}ab$:
\begin{align}
[ab,c]&=[a,c]^{b}\,[b,c], \label{eq:id1}\\
[a,bc]&=[a,c]\,[a,b]^{c}, \label{eq:id2}\\
[a,b^{-1}]&=\big([a,b]^{-1}\big)^{b^{-1}}, \label{eq:id3}\\
[a^{-1},b]&=\big([a,b]^{-1}\big)^{a^{-1}}, \label{eq:id4}\\
[a^{f},b]&=[a,\,b^{\,f^{-1}}]^{\,f}, \label{eq:id5}\\
[a,b]&=[b,a]^{-1}. \label{eq:id6}
\end{align}
Each is verified by direct expansion. (For \eqref{eq:id4}: put $a^{-1}$ for $a$ in
\eqref{eq:id1} to get $1=[a^{-1},c]^{a}[a,c]$. For \eqref{eq:id5}:
$[a^f,b]=f^{-1}a^{-1}f\,b^{-1}\,f^{-1}af\,b=[a,b^{f^{-1}}]^{f}$.) All six have been
verified by direct word reduction in the free group.

\begin{lemma}[Commutator engine]\label{lem:engine}
Let $S\subseteq F$ and $H=\langle S\rangle^{F}$. Then $H$ is normal and
\begin{equation}\label{eq:engine}
[H,F]=\big\langle\,[s,x] : s\in S,\ x\in X^{\pm}\,\big\rangle^{F}.
\end{equation}
\end{lemma}
\begin{proof}
$H$ is normal by construction. Write
$M=\langle\,[s,x]:s\in S,\,x\in X^{\pm}\,\rangle^{F}$. Each $[s,x]\in[H,F]$, which is
normal; hence $M\subseteq[H,F]$.

For the reverse inclusion, $[H,F]$ is generated by the $[h,y]$ with $h\in H$,
$y\in F$; as $M$ is normal it suffices to prove $[h,y]\in M$ for all such.

\emph{Step 1: $[s,y]\in M$ for $s\in S$, $y\in F$.} Induct on the word length
$\ell$ of $y$ in $X^{\pm}$. For $\ell=0$, $[s,1]=1$. For $\ell=1$: $y=x\in X$ gives
$[s,x]\in M$; $y=x^{-1}$ gives $[s,x^{-1}]=([s,x]^{-1})^{x^{-1}}\in M$ by
\eqref{eq:id3} and normality. For $\ell\ge2$ write $y=y'z$, $z\in X^{\pm}$; by
\eqref{eq:id2}, $[s,y]=[s,z]\,[s,y']^{z}\in M$ by induction and normality.

\emph{Step 2: $[s^{f},x]\in M$ for $s\in S$, $f\in F$, $x\in X^{\pm}$.} By
\eqref{eq:id5}, $[s^{f},x]=[s,x^{f^{-1}}]^{f}\in M$ by Step~1 and normality.

\emph{Step 3: $[t,y]\in M$ for $t\in\{(s^{f})^{\pm1}\}$, $y\in F$.} For $t=s^f$,
Step~2 plus the word induction of Step~1 (verbatim with $s^f$ for $s$) gives
$[s^f,y]\in M$. For $t=(s^f)^{-1}$, \eqref{eq:id4} gives
$[(s^f)^{-1},x]=([s^f,x]^{-1})^{(s^f)^{-1}}\in M$, and the same word induction gives
$[(s^f)^{-1},y]\in M$.

\emph{Step 4: $[h,y]\in M$ for $h\in H$.} Write $h=t_1\cdots t_n$ with each
$t_i\in\{(s_i^{f_i})^{\pm1}\}$ and induct on $n$ via \eqref{eq:id1}:
$[h,y]=[t_1,y]^{h'}[h',y]$, $h'=t_2\cdots t_n$; both factors lie in $M$ by Step~3,
the inductive hypothesis, and normality.
\end{proof}

\begin{lemma}\label{lem:caseA}
If $u\in N_d$ (any $d$) and $v\in F$, then $[u,v]\in N_d$ and $[v,u]\in N_d$.
\end{lemma}
\begin{proof}
$N_d$ is normal, so $u^{v}\in N_d$ and $[u,v]=u^{-1}u^{v}\in N_d$; also
$[v,u]=(u^{-1})^{v}u\in N_d$.
\end{proof}

\begin{lemma}\label{lem:easy}
For every $c\ge1$, $N_c\subseteq\gamma_c(F)$.
\end{lemma}
\begin{proof}
By Proposition~\ref{prop:hall} every $b\in B_c$ lies in the normal subgroup
$\gamma_c(F)$; hence $N_c=\langle B_c\rangle^{F}\subseteq\gamma_c(F)$.
\end{proof}

\section{Reducing the theorem to the single-level statement \eqref{eq:S}}
\label{sec:chain}

\begin{lemma}[Monotonicity from \eqref{eq:S}]\label{lem:mono}
If $B_w\subseteq N_{w-1}$ for some $w\ge2$, then $N_w\subseteq N_{w-1}$.
\end{lemma}
\begin{proof}
$N_{w-1}$ is a normal subgroup containing $B_w$, hence contains $N_w=\langle
B_w\rangle^{F}$.
\end{proof}

\begin{lemma}[$B_{\ge c}\subseteq N_c$ from \eqref{eq:S}]\label{lem:Bgec}
Assume \eqref{eq:S}. Then for every $c\ge1$ and every $w\ge c$, $B_w\subseteq N_c$;
hence $B_{\ge c}\subseteq N_c$.
\end{lemma}
\begin{proof}
Fix $c$ and induct on $w\ge c$. For $w=c$, $B_c\subseteq N_c$ by definition. For
$w>c$: by \eqref{eq:S} and Lemma~\ref{lem:mono}, $N_w\subseteq N_{w-1}$, so
$B_w\subseteq N_w\subseteq N_{w-1}$. By the inductive hypothesis (excess
$w-1-c<w-c$, $w-1\ge c$), $B_{w-1}\subseteq N_c$, hence
$N_{w-1}=\langle B_{w-1}\rangle^{F}\subseteq N_c$. Therefore $B_w\subseteq
N_{w-1}\subseteq N_c$.
\end{proof}

\begin{proposition}[Chaining reduction]\label{prop:chain}
The following are equivalent:
\begin{enumerate}
\item[\textup{(a)}] $\gamma_c(F)=N_c$ for every $c\ge1$ \textup{(Theorem~\ref{thm:main})};
\item[\textup{(b)}] $B_{\ge c}\subseteq N_c$ for every $c\ge1$;
\item[\textup{(c)}] the single-level statement \eqref{eq:S}: $B_w\subseteq N_{w-1}$ for every $w\ge2$.
\end{enumerate}
\end{proposition}
\begin{proof}
(a)$\Rightarrow$(b): if $\gamma_c(F)=N_c$ then each $b\in B_w$ ($w\ge c$) lies in
$\gamma_w(F)\subseteq\gamma_c(F)=N_c$ (Proposition~\ref{prop:hall}).

(b)$\Rightarrow$(c): take $c=w-1$; then $B_w\subseteq B_{\ge w-1}\subseteq N_{w-1}$.

(c)$\Rightarrow$(b) is Lemma~\ref{lem:Bgec}.

(b)$\Rightarrow$(a): By Proposition~\ref{prop:hall}, $\gamma_c(F)=\langle B_{\ge
c}\rangle$ as a subgroup; by (b), $B_{\ge c}\subseteq N_c$, so
$\gamma_c(F)\subseteq N_c$. With Lemma~\ref{lem:easy}, $\gamma_c(F)=N_c$.
\end{proof}

By Proposition~\ref{prop:chain} it suffices to prove the single-level statement
\eqref{eq:S}. This is the precise, finitary form of the issue's demand: instead of
``every basic commutator of weight $w>c$ lies in $N_c$ for every $(w,c)$,'' we must
prove only that each weight-$w$ basic commutator lies in the normal closure of the
weight-$(w-1)$ basic commutators.

\section{The single-level statement \eqref{eq:S} and its hard core}\label{sec:hard}

Fix $w\ge2$ and let $b=[u,v]\in B_w$, so $u,v$ are basic commutators with $u>v$,
$\mathrm{wt}(u)=p$, $\mathrm{wt}(v)=q$, $p+q=w$, and (since the Hall order refines
weight and $u>v$) $p\ge q\ge1$, with $p,q<w$.

\begin{lemma}[Reduction of \eqref{eq:S} to single commutators]\label{lem:S-to-engine}
If, for the given $w$,
\begin{equation}\label{eq:dx}
[d,x]\in N_{w-1}\qquad\text{for every basic commutator $d$ with $\mathrm{wt}(d)\ge w-1$ and every $x\in X^{\pm}$,}
\end{equation}
then $B_w\subseteq N_{w-1}$.
\end{lemma}
\begin{proof}
By Proposition~\ref{prop:hall}, $\gamma_{w-1}(F)=\langle B_{\ge w-1}\rangle$ as a
subgroup; being normal, $\gamma_{w-1}(F)=\langle B_{\ge w-1}\rangle^{F}$. By the
engine (Lemma~\ref{lem:engine}) with $S=B_{\ge w-1}$,
\[
\gamma_w(F)=[\gamma_{w-1}(F),F]=\big\langle\,[d,x] : d\in B_{\ge w-1},\ x\in
X^{\pm}\,\big\rangle^{F}.
\]
By \eqref{eq:dx} every generator $[d,x]$ lies in the normal subgroup $N_{w-1}$, so
$\gamma_w(F)\subseteq N_{w-1}$, and $B_w\subseteq\gamma_w(F)\subseteq N_{w-1}$.
\end{proof}

We now analyse \eqref{eq:dx} by the weight of $d$.

\begin{lemma}[The case $\mathrm{wt}(d)=w-1$]\label{lem:dw1}
If $d\in B_{w-1}$ and $x\in X^{\pm}$, then $[d,x]\in N_{w-1}$.
\end{lemma}
\begin{proof}
$d\in B_{w-1}\subseteq N_{w-1}$ by definition; $N_{w-1}$ is normal, so
$[d,x]\in N_{w-1}$ by Lemma~\ref{lem:caseA}.
\end{proof}

The remaining case of \eqref{eq:dx} is $\mathrm{wt}(d)=m\ge w$, equivalently the
membership $[d,x]\in N_{w-1}$ for $d\in B_m$, $m\ge w$. Since $d\in B_m\subseteq
B_{\ge w-1}$, this is implied by $B_{\ge w-1}\subseteq N_{w-1}$, i.e.\ by
Theorem~\ref{thm:nc} at $c=w-1$. Thus:

\begin{lemma}[Discharge of \eqref{eq:dx}]\label{lem:hardcase}
For every $w\ge2$, \eqref{eq:dx} holds; hence (Lemma~\ref{lem:S-to-engine})
$B_w\subseteq N_{w-1}$, i.e.\ \eqref{eq:S} holds.
\end{lemma}
\begin{proof}
Let $d$ be a basic commutator with $\mathrm{wt}(d)=m\ge w-1$ and $x\in X^{\pm}$. By
Theorem~\ref{thm:nc} applied at $c=w-1$, $B_{\ge w-1}\subseteq N_{w-1}$; since
$d\in B_m\subseteq B_{\ge w-1}$, we have $d\in N_{w-1}$. As $N_{w-1}$ is normal,
$[d,x]\in N_{w-1}$ by Lemma~\ref{lem:caseA}. This is \eqref{eq:dx}.
\end{proof}

Before assembling the proof we make explicit how the issue's commutators occur,
and why the elementary part stops exactly at \eqref{eq:dx} with $\mathrm{wt}(d)\ge
w$.

\begin{lemma}[Peeling a commutator second entry]\label{lem:peel}
For all $a,e_1,e_2\in F$,
\begin{equation}\label{eq:peel}
[a,[e_1,e_2]]=[a,e_2]\,[a,e_1]^{e_2}\,[a,e_2^{-1}]^{e_1e_2}\,[a,e_1^{-1}]^{e_2^{-1}e_1e_2}.
\end{equation}
Hence $[a,[e_1,e_2]]$ lies in the normal closure of $\{[a,e_1],[a,e_2]\}$ in $F$.
\end{lemma}
\begin{proof}
Write $[e_1,e_2]=e_1^{-1}e_2^{-1}e_1e_2$ and peel the rightmost letter repeatedly
using \eqref{eq:id2}:
\begin{align*}
[a,e_1^{-1}e_2^{-1}e_1e_2]
&=[a,e_2]\,[a,e_1^{-1}e_2^{-1}e_1]^{e_2}\\
&=[a,e_2]\,\big([a,e_1]\,[a,e_1^{-1}e_2^{-1}]^{e_1}\big)^{e_2}\\
&=[a,e_2]\,[a,e_1]^{e_2}\,\big([a,e_2^{-1}]\,[a,e_1^{-1}]^{e_2^{-1}}\big)^{e_1e_2}\\
&=[a,e_2]\,[a,e_1]^{e_2}\,[a,e_2^{-1}]^{e_1e_2}\,[a,e_1^{-1}]^{e_2^{-1}e_1e_2},
\end{align*}
which is \eqref{eq:peel}. By \eqref{eq:id3}, $[a,e_i^{-1}]$ is a conjugate of
$[a,e_i]^{-1}$, so every factor of \eqref{eq:peel} is a conjugate of
$[a,e_1]^{\pm1}$ or $[a,e_2]^{\pm1}$.
\end{proof}

\begin{remark}[Why elementary descent does not finish]\label{rem:peel}
Lemma~\ref{lem:peel} writes the issue's commutator $[[x_2,x_1],[x_3,x_1]]$ (weight
$4$, $c=3$) explicitly as a product of conjugates of $[[x_2,x_1],x_3]$ and
$[[x_2,x_1],x_1]$, two left-normed commutators of weight $3$. Iterating
\eqref{eq:peel} reduces any $[u,v]$ to commutators $[u,x]$ with $x\in X^{\pm}$, of
weight $\mathrm{wt}(u)+1\le w-1$. These have weight \emph{below} $w$ and, when of
weight $<w-1$, cannot lie in $N_{w-1}$ for weight reasons (an element of $N_{w-1}$
lies in $\gamma_{w-1}(F)$, by Lemma~\ref{lem:easy}). Thus peeling \emph{lowers}
weight and cannot by itself place a weight-$w$ basic commutator in $N_{w-1}$. The
same phenomenon afflicts every elementary route: by Lemma~\ref{lem:moddescent} the
graded calculus controls $\gamma_w(F)$ only modulo $\gamma_{w+1}(F)$, leaving the
infinite tail $\bigcap_m N_{w-1}\gamma_m(F)$, whose collapse to $N_{w-1}$ is
equivalent to the conclusion (\S\ref{sec:why}). This is exactly the residue
$\mathrm{wt}(d)\ge w$ in \eqref{eq:dx}, and it is supplied by Theorem~\ref{thm:nc}.
\end{remark}

\begin{proof}[Proof of Theorem~\ref{thm:main}]
By Lemma~\ref{lem:hardcase}, the single-level statement \eqref{eq:S} holds; by
Proposition~\ref{prop:chain}, $\gamma_c(F)=N_c$ for every $c\ge1$.
\end{proof}

\begin{remark}[Logical content]
The chain is: Theorem~\ref{thm:nc} (classical, at level $w-1$) $\Rightarrow$
\eqref{eq:dx} (Lemma~\ref{lem:hardcase}) $\Rightarrow$ \eqref{eq:S}
(Lemma~\ref{lem:S-to-engine}) $\Rightarrow$ Theorem~\ref{thm:main}
(Proposition~\ref{prop:chain}). Of course Theorem~\ref{thm:nc} \emph{is}
Theorem~\ref{thm:main}; the value of the elementary part is the equivalence
\textup{(a)}$\Leftrightarrow$\textup{(c)} of Proposition~\ref{prop:chain}, which
shows that to prove the theorem for all $(w,c)$ it suffices to prove the single
weight-by-one statement \eqref{eq:S}, and that all of \eqref{eq:dx} except the case
$\mathrm{wt}(d)\ge w$ is elementary. We are transparent that the genuinely hard
content -- the case $\mathrm{wt}(d)\ge w$ -- is classical and equal to the theorem
one level up, and we cite it rather than disguise it.
\end{remark}

\begin{corollary}
For every $c\ge1$, $B_{\ge c}\subseteq N_c$ and $F/N_c$ is nilpotent of class
$\le c-1$.
\end{corollary}
\begin{proof}
$B_{\ge c}\subseteq\gamma_c(F)=N_c$ by Theorem~\ref{thm:main} and
Proposition~\ref{prop:chain}(b); $\gamma_c(F/N_c)=\gamma_c(F)N_c/N_c=1$.
\end{proof}

\section{Why a genuine classical input is necessary, and non-circularity}
\label{sec:why}

\subsection{Graded data plus elementary calculus do not suffice}

\begin{lemma}[Mod-$\gamma$ descent]\label{lem:moddescent}
For every $w\ge c\ge1$, $\gamma_w(F)\subseteq N_c\,\gamma_{w+1}(F)$.
\end{lemma}
\begin{proof}
Induct on $w\ge c$. For $w=c$, \eqref{eq:hall-decomp} gives
$\gamma_c(F)=\langle B_c\rangle\gamma_{c+1}(F)\subseteq N_c\gamma_{c+1}(F)$. For
$w>c$ assume $\gamma_{w-1}(F)\subseteq N_c\gamma_w(F)$. By
Proposition~\ref{prop:hall}, $\gamma_{w-1}(F)=\langle B_{\ge w-1}\rangle^F$, so by
the engine $\gamma_w(F)$ is the normal closure of the $[a,x]$, $a\in B_{\ge w-1}$,
$x\in X^{\pm}$. Writing $a=nt$ with $n\in N_c$, $t\in\gamma_w(F)$, \eqref{eq:id1}
gives $[a,x]=[n,x]^t[t,x]$ with $[n,x]\in N_c$ (Lemma~\ref{lem:caseA}) and
$[t,x]\in\gamma_{w+1}(F)$, so $[a,x]\in N_c\gamma_{w+1}(F)$, a normal subgroup.
Hence $\gamma_w(F)\subseteq N_c\gamma_{w+1}(F)$.
\end{proof}

\begin{lemma}[Identical graded data]\label{lem:gradedNc}
For every $w\ge1$ the images of $N_c$ and of $\gamma_c(F)$ in
$L_w:=\gamma_w(F)/\gamma_{w+1}(F)$ coincide; equivalently $N_c$ and $\gamma_c(F)$
have the same image in every nilpotent quotient $F/\gamma_m(F)$. This holds
independently of Theorem~\ref{thm:main}.
\end{lemma}
\begin{proof}
For $w\ge c$, Lemma~\ref{lem:moddescent} and Dedekind's modular law give
$\gamma_w(F)=(N_c\cap\gamma_w(F))\gamma_{w+1}(F)$, so $N_c$ surjects onto $L_w$, as
does $\gamma_c(F)\supseteq\gamma_w(F)$. For $1\le w<c$, both subgroups lie in
$\gamma_{w+1}(F)$ with trivial image in $L_w$.
\end{proof}

\begin{example}\label{ex:Z}
Equal graded data does not force equality, even with residual nilpotence. In
$G=\mathbb Z$ with filtration $F_w=2^{w-1}\mathbb Z$ one has $[F_i,F_j]=0$,
$\bigcap_w F_w=0$, $F_w/F_{w+1}\cong\mathbb Z/2$. For $A=3\mathbb Z\subsetneq\mathbb
Z=B$, both $A\cap F_w$ and $B\cap F_w$ surject onto $\mathbb Z/2$ for all $w$, so
$A,B$ have identical graded data yet $A\ne B$.
\end{example}

\begin{remark}\label{rem:graded-insufficient}
By Lemma~\ref{lem:gradedNc} no invariant of the associated graded data of $F$ can
separate $N_c$ from $\gamma_c(F)$, and Example~\ref{ex:Z} shows such data are in
principle insufficient to force equality of subgroups. The missing fact is
$\bigcap_m N_c\gamma_m(F)=N_c$, i.e.\ residual nilpotence of $F/N_c$, which is not a
graded invariant of $F$. Hence the equality $\gamma_c(F)=N_c$ cannot be deduced from
Proposition~\ref{prop:hall} (the graded basis theorem) together with the elementary
commutator calculus of \S\ref{sec:engine} alone; a genuine input of the strength of
Theorem~\ref{thm:nc} is required. This justifies, rather than hides, our single
classical citation.
\end{remark}

\subsection{The cited input is not the conclusion in disguise}

\emph{(a) What is proved here.} Everything except Theorem~\ref{thm:nc} is proved in
this paper from the graded basis theorem (Proposition~\ref{prop:hall}) and
elementary commutator calculus: the engine (Lemma~\ref{lem:engine}), the chaining
reduction (Proposition~\ref{prop:chain}), the reduction to single commutators
(Lemma~\ref{lem:S-to-engine}), the peeling identity (Lemma~\ref{lem:peel}), and the
mod-$\gamma$ descent (Lemma~\ref{lem:moddescent}). In particular the equivalence of
Theorem~\ref{thm:main} with the single-level statement \eqref{eq:S} is new and
elementary.

\emph{(b) The previous circular citation is removed.} The earlier version cited the
equality $\gamma_w=N_w$ under the name ``closure identity of the collecting
process'' and applied it at level $c+1$ to finish; that is the conclusion restated.
We have removed it. We instead cite Theorem~\ref{thm:nc} -- the normal-generation
form of the basis theorem -- and are explicit that it equals the conclusion. We do
not claim to reduce it to anything weaker; by Remark~\ref{rem:graded-insufficient}
no such reduction from graded data is possible.

\emph{(c) The cited theorem's classical proof is non-circular.} The standard proofs
of Theorem~\ref{thm:nc} (Hall; Magnus--Karrass--Solitar) proceed by the collecting
process, a terminating word-rewriting algorithm using only
\eqref{eq:id1}--\eqref{eq:id2}, whose termination is by a weight-lexicographic
measure on the multiset of uncollected entries of a finite word. That algorithm
never mentions the normal closures $N_c$, residual nilpotence of $F/N_c$, or
Theorem~\ref{thm:main}. Hence importing Theorem~\ref{thm:nc} as a classical result
does not presuppose the present development.
\end{document}
